\section{Vanishing of the limiting energy in two dimensions}
\label{sec:vanishing}

Throughout this section, the disorder strength and the interior energy are
fixed. All auxiliary constructions are those of the preceding sections.
In particular, \cref{screen:limit} provides finite numbers $d_n$ with
$d_n\to d\geq0$. No convergence rate uniform in $n$ is assumed.

\begin{proposition}\label{van:zero}
The limiting affine energy vanishes: $d=0$.
\end{proposition}

Suppose for contradiction that $d>0$. We first thin the gaps between
retained inner squares and grid walls. A localized change of the affine
data disappears after the complementary thinning pass. Since the
expected initial and final affine energy densities approach the same limit,
choosing the sign of this change forces its local first variation to be
small in expectation.

We then place weak probes in a square annulus between one inner square
and its wall. A logarithmic voltage profile makes total transmission to
the probes small, so the cell-to-cell scattering block is nearly
unitary. Its surviving internal jump energy gives an analytic trace a
nonzero central value. Closing the probes identifies the boundary
modulus with the preceding first variation evaluated after a
support-preserving shift of the true potentials. A bounded cost for
changing the law transfers smallness to boundary probabilities, and
dimension-independent moment bounds turn this into a small boundary
integral, contradicting submean.

\subsection{Geometry and order of scales}

Fix a number $M_0>20$, to be enlarged below. After choosing $n$, choose
$m>n$ and a power of three $L_*$ such that
\begin{equation}\label{van:scales}
 e^{b_m^2}\leq L_*<3e^{b_m^2},
 \qquad R=\frac{L_*}{M_0^2}.
\end{equation}
The torus side $b_N$ will be much larger than $L_*$. Since $L_*$ is a
power of three, the torus admits an aligned partition into grid tiles
$D$ of side $L_*$. The alignment is specified by the hierarchy offset.
Write $c_D$ for a tile center and use max-norm radii about this center.

Let $I_D$ be the union of level-$m$ tiles whose centers have radius at
most $R$, and let $O_D$ be the union of level-$m$ tiles whose centers
have radius at least $L_*/3$. Thus the inner square has radius
$R+O(b_m)$, and the outer walls begin at radius $L_*/3+O(b_m)$. The
remaining level-$m$ tiles form the gap. Thin levels $n,\ldots,m-1$ in
all gaps, leaving the inner squares and walls at stage $n$. The result
is called the intermediate system. A complementary pass on the retained
regions finishes the transition to uniform stage $m$. Both pass regions
are unions of whole level-$m$ tiles, as required by
\cref{hier:comparison}.

Define the selector event
\begin{equation}\label{van:no-large}
 \mathcal P_D
 =\{\text{no selected tile of level at least $m$ meets $D$}\}.
\end{equation}
This event involves selectors only. Conditional on the offset,
\begin{equation}\label{van:no-large-bound}
 \PP(\mathcal P_D^c\mid\text{offset})
 \leq C\sum_{l\geq m}(1+L_*/b_l)^2q_l=o(1).
\end{equation}
Indeed $q_l=e^{-b_l^6}$ and $L_*<3e^{b_m^2}$, so the right side tends
to zero faster than any fixed negative power of $R$. On $\mathcal P_D$,
the intermediate cell ports in $D$ lie only in $I_D\cup O_D$; after the
complementary pass there are no cell ports in $D$.

Independently of the potentials, selectors, and acceptance variables,
choose a uniformly random grid tile, conditional on the offset. For
this test tile write $D,c,I,O,\mathcal P$ instead of
$D,c_D,I_D,O_D,\mathcal P_D$. All expectations below include this
choice. It is a permitted independent prescription for the screening
and comparison estimates.
The retained regions and the probe annulus used below are shown in
\cref{van:geometry-figure}.

\begin{figure}[t]
\centering
\resizebox{0.92\linewidth}{!}{%
\begin{tikzpicture}[x=1cm,y=1cm,>=Stealth, line width=0.55pt]
  \fill[black!12,even odd rule]
    (-3.25,-3.25) rectangle (3.25,3.25)
    (-2.60,-2.60) rectangle (2.60,2.60);
  \fill[blue!14,even odd rule]
    (-1.94,-1.94) rectangle (1.94,1.94)
    (-1.45,-1.45) rectangle (1.45,1.45);
  \fill[green!12] (-0.70,-0.70) rectangle (0.70,0.70);
  \draw[thick] (-3.25,-3.25) rectangle (3.25,3.25);
  \draw (-2.60,-2.60) rectangle (2.60,2.60);
  \draw[blue!65!black] (-1.94,-1.94) rectangle (1.94,1.94);
  \draw[blue!65!black] (-1.45,-1.45) rectangle (1.45,1.45);
  \draw[green!45!black] (-0.70,-0.70) rectangle (0.70,0.70);
  \draw[dashed,black!50] (-0.96,-0.96) rectangle (0.96,0.96);
  \node at (0,0.18) {$I_D$};
  \node[font=\scriptsize] at (0,-0.16) {retained};
  \node[font=\small] at (0,2.27) {thinned gap};
  \node[font=\small] at (0,1.69) {probe annulus $C$};
  \node[font=\small] at (0,2.94) {retained wall $O_D$};
  \fill (0,0) circle (0.025);
  \draw[->] (0,-0.35) -- (3.52,-0.35);
  \node[anchor=west,font=\scriptsize] at (3.56,-0.35) {radius};
  \draw (0.70,-0.43) -- (0.70,-0.27);
  \draw (1.45,-0.43) -- (1.45,-0.27);
  \draw (1.94,-0.43) -- (1.94,-0.27);
  \draw (2.60,-0.43) -- (2.60,-0.27);
  \draw (3.25,-0.43) -- (3.25,-0.27);
  \draw[black!55] (0.70,-0.44) -- (0.70,-0.85) -- (4.05,-0.85);
  \draw[black!55] (1.45,-0.44) -- (1.45,-1.38) -- (4.05,-1.38);
  \draw[black!55] (1.94,-0.44) -- (1.94,-1.91) -- (4.05,-1.91);
  \draw[black!55] (2.60,-0.44) -- (2.60,-2.44) -- (4.05,-2.44);
  \draw[black!55] (3.25,-0.44) -- (3.25,-2.97) -- (4.05,-2.97);
  \node[anchor=west,font=\small] at (4.12,-0.85) {$R+O(b_m)$};
  \node[anchor=west,font=\small] at (4.12,-1.38) {$M_0R$};
  \node[anchor=west,font=\small] at (4.12,-1.91) {$2M_0R$};
  \node[anchor=west,font=\small] at (4.12,-2.44) {$L_*/3+O(b_m)$};
  \node[anchor=west,font=\small] at (4.12,-2.97) {$L_*/2$};
  \node[align=left,anchor=west,font=\small] at (3.95,1.05)
    {$L_*=M_0^2R$\\$b_m\ll R$\\$M_0>20$};
  \node[font=\scriptsize,align=center] at (0,-3.63)
    {Dashed square: adjustment support ends at radius $2R$.};
\end{tikzpicture}%
}
\caption{One grid tile in the intermediate system. The annulus is a
square annulus, because all indicated radii use the max norm. The
picture is schematic, not to scale: in the proof $M_0$ is large, so
$R\ll M_0R<2M_0R\ll L_*/3$. Inner and wall boundaries are rounded to
unions of level-$m$ tiles; the probe annulus is not so rounded. On
$\mathcal P_D$, the gap has no cell ports, and the added probes occur
only in the displayed annulus.}
\label{van:geometry-figure}
\end{figure}

\subsection{An adaptive adjustment of the affine data}

For the moment use no artificial probes. Let $S'$ be the intermediate
scattering matrix and let $f=f_N$ be the four-component torus embedding
used to define the affine energy in \cref{screen:limit}. Its chord
distances are globally comparable to shortest Euclidean torus
distances, and its nearest-neighbor increments are bounded. Define
four-component data on the intermediate cell ports by
\begin{equation}\label{van:trace-data}
 F=\frac{f-f(c)}{R},\qquad F_I=\one_I F,
 \qquad
 H^0=\sum_{\nu=1}^{4}
 \tr\!\left[F_{I,\nu}(F_\nu-S'F_\nu S'^*)\right].
\end{equation}
As elsewhere, real scalar data inside a matrix expression denote the
corresponding diagonal matrices. The trace $H^0$ is real, since every
diagonal entry of $S'F_\nu S'^*$ is real.
It is the local first variation of the affine energy: for fixed $S'$,
\[
 \left.\frac{\dd}{\dd t}\mathcal Q_{S'}(f-tRF_I)\right|_{t=0}
       =-2R^2H^0.
\]
Here the perturbation is data on the intermediate ports. The proof below
extends it by a cutoff to site data with bounded increments and chooses
its sign separately in each grid tile. On the selector event, the
affected ports disappear in the complementary pass.

\begin{lemma}\label{van:adjustment}
Fix $M_0$. For every $a>0$, one can arrange
\begin{equation}\label{van:adjustment-small}
 \E\bigl[\one_{\mathcal P}|H^0|\bigr]\leq a
\end{equation}
by taking $n$, then $m>n$, then $N$ sufficiently large. More precisely,
the estimate holds for all sufficiently large $n$, all sufficiently
large $m$ depending on $n$, and all sufficiently large $N$ depending on
the fixed pair $(n,m)$.
\end{lemma}

\begin{proof}
For each grid tile choose a cutoff $\eta_D$ that is one on $I_D$,
vanishes at radius $2R$, and has Lipschitz constant at most $C/R$.
Such cutoffs exist once $b_m/R$ is small. Define $H_D^0$ by
\cref{van:trace-data} with $D$ as the test tile, and set
\[
 \varphi(x)=-\sum_D
   \one_{\mathcal P_D}\operatorname{sgn}(H_D^0)
   \eta_D(x)\bigl(f(x)-f(c_D)\bigr),
 \qquad \operatorname{sgn}(0)=0.
\]
The supports are disjoint. On each support, $|f-f(c_D)|\leq CR$,
so the product has uniformly bounded nearest-neighbor increments.
Across support boundaries the same bound holds because the cutoff
vanishes there. Using a shortest lattice path and global comparability
of chord distances gives
\begin{equation}\label{van:adjustment-lipschitz}
 |\varphi(x)-\varphi(y)|\leq C|f(x)-f(y)|
 \quad\text{for all torus sites }x,y.
\end{equation}

This random adjustment is admissible at the start of the complementary
pass in \cref{hier:comparison}. The selector events may be used as
pre-information. Also $S'$ uses the calibrated reference value at every
still-open root, not the hidden true value at that root. By
\cref{hier:hidden}, the disjoint eligible cells and their recorded
statuses do not expose any of those hidden values. Thus $H_D^0$ and
its sign are measurable with respect to the allowed pre-information.
The adjustment does not use the subsequent thinning marks.

On $\mathcal P_D$, all intermediate ports meeting $\supp\eta_D$ lie
in $I_D$, where $\eta_D=1$. No final port meets this support. Hence
$g=f+t\varphi$, $0<t<1$, agrees with $f$ on the final stage-$m$
ports. By symmetry and unitarity of $S'$, the quadratic energy expands as
\[
 \mathcal Q_{S'}(g)=\mathcal Q_{S'}(f)
 +2t\sum_\nu\tr\!\left[
     \varphi_\nu(f_\nu-S'f_\nu S'^*)\right]
 +t^2\mathcal Q_{S'}(\varphi).
\]
One may subtract $f_\nu(c_D)$ in the trace belonging to a fixed tile,
because $S'$ is unitary. Its linear contribution is therefore
$-2tR^2\one_{\mathcal P_D}|H_D^0|$. There are $b_N^2/L_*^2$ grid
tiles. Averaging the uniformly selected test tile and dividing by the
torus area yields
\begin{equation}\label{van:adjustment-expansion}
 \frac{\E\mathcal Q_{S'}(g)}{b_N^2}
 =\frac{\E\mathcal Q_{S'}(f)}{b_N^2}
 -2tM_0^{-4}\E[\one_{\mathcal P}|H^0|]
 +t^2\frac{\E\mathcal Q_{S'}(\varphi)}{b_N^2}.
\end{equation}
By \cref{van:adjustment-lipschitz}, the last energy is at most a
constant times $\mathcal Q_{S'}(f)$.

Apply \cref{hier:comparison,hier:charges} to the first pass with data
$f$. Then apply the comparison to the complementary pass with data
$g$, using the pointwise increment bound for $g$ obtained from
\cref{van:adjustment-lipschitz}. The allowed increment constant is
uniform for $0<t<1$. The fixed-stage affine limits give
\begin{align*}
 b_N^{-2}\E\mathcal Q_{S'}(f)
   &\leq d_n+\varepsilon_n+o_{n,m;N}(1),\\
 b_N^{-2}\E\mathcal Q_{S'}(g)
   &\geq d_m-\varepsilon_n-o_{n,m;N}(1),
\end{align*}
where $\varepsilon_n\to0$ uniformly in $m>n$. Here and below,
$o_{n,m;N}(1)$ tends to zero as $N\to\infty$ with $n,m$ fixed.
Boundedness of $d_n$ now gives, from
\cref{van:adjustment-expansion},
\begin{equation}\label{van:adjustment-inequality}
 2tM_0^{-4}\E[\one_{\mathcal P}|H^0|]
 \leq Ct^2+d_n-d_m+2\varepsilon_n+o_{n,m;N}(1).
\end{equation}
First take $t$ sufficiently small that the quadratic contribution,
after division by $2tM_0^{-4}$, is less than $a/2$. Since $d_n$ is
convergent, the remaining terms are less than the required tolerance
for sufficiently large $n,m,N$ in the stated order. This proves
\cref{van:adjustment-small}.
\end{proof}

\subsection{An annulus of artificial probes}

Add probes at all sites of
\[
 C=\{x:M_0R\leq |x-c|_\infty\leq 2M_0R\},
\]
and nowhere else, using strength $\tau\vartheta(v_x)$, where
\begin{equation}\label{van:probe-scale}
 \tau=\frac{C_0}{R}.
\end{equation}
The constant $C_0$ will be fixed before $n,m,N$ are chosen. For
sufficiently large $m$, \cref{van:scales} implies
$\tau\leq e^{-c_*b_m}$, so this mask is allowed throughout the first
comparison pass. Use the same underlying hierarchy choices as before.
The resulting intermediate scattering matrix has block form
\begin{equation}\label{van:block}
 S=\begin{pmatrix}A&T\\T^\top&B\end{pmatrix}.
\end{equation}
The first block consists of all cell ports, and the second consists of
the new probes. We also use $C$ for their index set. In this section
$I$ denotes the cell ports in the retained inner square and $J$ denotes
all other cell ports; these symbols are local to this section. On
$\mathcal P$, every port in $J$ is outside the probe annulus toward
the exterior. Set
\[
 K=\norm{T}_{\mathrm{HS}}^2.
\]
The data $F,F_I$ still refer to \cref{van:trace-data}, now restricted
to the cell ports of \cref{van:block}.

We first give the deterministic use of a good annular barrier. A site
is called good for this purpose when $\vartheta(v_x)\geq d_0$, where
$d_0>0$ is fixed sufficiently small. Count only edges with both
endpoints good and lying in a buffered square band strictly inside
$C$, for example with both endpoint radii in
$[1.2M_0R,1.8M_0R]$. Give such an edge cost one, and all other edges
cost zero. Denote by $\ell(x)$ the resulting cost distance from the
region $|x-c|_\infty\geq 2M_0R$. Let $\mathcal B$ be the event
\begin{equation}\label{van:barrier-event}
 \ell(x)\geq\gamma M_0R
 \quad\text{whenever }|x-c|_\infty\leq M_0R,
\end{equation}
where $\gamma>0$ is a sufficiently small fixed constant.

\begin{lemma}\label{van:barrier}
For suitable fixed $d_0,\gamma>0$,
$\PP(\mathcal B^c)\leq C e^{-cM_0R}$. On
$\mathcal B\cap\mathcal P$, for every compatible matrix $D'$ on the
source columns indexed by $J$,
\begin{equation}\label{van:barrier-bound}
 \norm{A_{IJ}D'}_{\mathrm{HS}}
 \leq \frac{C}{M_0R\tau d_0}
       \norm{S_{CJ}D'}_{\mathrm{HS}}.
\end{equation}
The constant includes the fixed factor $\gamma^{-1}$. The same cutoff
argument controls every absorbed output supported where the cutoff
equals one; it requires no norm bound on the exterior resolvent.
\end{lemma}

\begin{proof}
Since $\vartheta$ is positive in the interior of the single-site
support, the iid probability
$p_0=\PP\{\vartheta(v_x)<d_0\}$ tends to zero with $d_0$. If
\cref{van:barrier-event} fails, trim a cheap crossing to the buffered
band and erase its loops. The resulting simple path has length
$j\geq cM_0R$. Choosing $\gamma$ sufficiently small, a fixed positive
fraction of its edges must touch bad sites. Since the path is simple,
a fixed positive fraction of its distinct vertices are bad. There are
at most $C(M_0R)^2 4^j$ candidate paths of length $j$, and at most
$2^{j+1}$ ways to designate a subset of their vertices. Independence
therefore bounds the failure probability by
\[
 C(M_0R)^2\sum_{j\geq cM_0R}4^j2^{j+1}p_0^{c'j}.
\]
For sufficiently small fixed $d_0$, this is exponentially small.

For the deterministic estimate, use the physical inverse $G$ of the
intermediate system with the probes inserted, and set
$u=GY_JD'$. Define
\[
 \chi(x)=\min\{1,\ell(x)/(\gamma M_0R)\}.
\]
It is zero at the $J$ sources and one at $I$. An edge across which
$\chi$ changes is a counted edge; on that edge
$|\chi(x)-\chi(y)|\leq(\gamma M_0R)^{-1}$. Since diagonal terms
commute with $\chi$, the equation for $\chi u$ has only the residual
\[
 q=[G^{-1},\chi]u
\]
and no original source. This residual is supported on the probe sites
at endpoints of changing edges. If $\sigma_x=\tau\vartheta(v_x)$,
write $q=Y_C a$ there. In each matrix entry of the map
$Y_C^\top u\mapsto a$, the conversion introduces
$\sigma_x^{-1/2}\sigma_y^{-1/2}\leq(\tau d_0)^{-1}$.
The vertex degree is bounded, and hence
\begin{equation}\label{van:commutator-source}
 \norm{a}_{\mathrm{HS}}
 \leq \frac{C}{M_0R\tau d_0}
       \norm{Y_C^\top u}_{\mathrm{HS}}.
\end{equation}
By \cref{mat:ward}, the norm of propagation from probe sources to
all absorbed outputs is at most one. Since $\chi=1$ at $I$,
\[
 \norm{Y_I^\top u}_{\mathrm{HS}}
 =\norm{Y_I^\top GY_Ca}_{\mathrm{HS}}
 \leq\norm{a}_{\mathrm{HS}}.
\]
The blocks in \cref{van:barrier-bound} are off diagonal, so replacing
the physical fields by their scattering blocks multiplies both sides
by the same factor $2$. This proves the assertion and also the stated
all-output version of the cutoff argument.
\end{proof}

The next event combines small total transmission with positive internal
jump energy and a bound on transmission weighted by $F$. The weighted
bound is needed because $F$ grows on distant cell ports: small $K$ alone
does not control those coordinate factors in the trace.

\begin{lemma}\label{van:good}
Assume $d>0$. There is a constant $c_1>0$ with the following property.
For every sufficiently small fixed $\kappa>0$, one can choose $M_0$
sufficiently large and finite constants $C_1,C_F$, before choosing
$C_0,n,m,N$, so that the following holds. For any fixed sufficiently
large $C_0$, and then $n,m,N$ sufficiently large successively, an event
$\mathcal G$ of probability at least $1/2$ satisfies
\begin{gather}
 \mathcal G\subset\mathcal P\cap\mathcal B,\qquad
 K\leq\kappa,\label{van:good-transmission}\\
 c_1d\leq
 \sum_{i,j\in I}|F(i)-F(j)|^2|A_{ij}|^2
 \leq C_1,\label{van:good-jumps}\\
 \norm{F_\nu T}_{\mathrm{HS}}\leq C_F
 \quad (1\leq\nu\leq4).\label{van:good-weighted}
\end{gather}
In particular, on this event,
\begin{equation}\label{van:external-small}
 \norm{A_{IJ}}_{\mathrm{HS}}
 +\sum_{\nu=1}^4\norm{A_{IJ}F_\nu|_J}_{\mathrm{HS}}
 \leq \frac{C(\sqrt\kappa+C_F)}{M_0C_0d_0}.
\end{equation}
The right side tends to zero as $C_0$ increases, with all constants
chosen before it kept fixed.
\end{lemma}

\begin{proof}
Fix small exception budgets for each of the following assertions, with
their sum less than $1/2$. All probability bounds are unconditional;
in particular, none of the screening or spatial-average estimates is
conditioned on $\mathcal P$. Let $D_*$ be a finite upper bound for all
$d_n$ after increasing the starting level if necessary.

\emph{Small total transmission to the annulus.}
Choose voltage one on $C$, zero within radius $2R$ and at radii at
least $L_*/4$, with logarithmic ramps in the intervening regions.
For example, as a function of $r=|x-c|_\infty$, the inner ramp is
$\log(r/(2R))/\log(M_0/2)$ and the outer ramp is
$\log((L_*/4)/r)/\log(M_0/8)$. It has a compact profile at scale $R$
and
\begin{equation}\label{van:capacity}
 \int_{\R^2}|\nabla\varphi|^2\,\dd x
 \leq\frac{C}{\log M_0}.
\end{equation}
Indeed the area element for square shells is a constant times
$r\,\dd r$, while the gradient on either ramp is at most
$C/(r\log M_0)$. On $\mathcal P$, all cell-port voltages are zero,
and all probe voltages are one. The intermediate energy is therefore
exactly $K$: each cell-to-probe pair is counted once by the symmetric
energy.

Use \cref{screen:profiles} for the starting stage-$n$ energy and
\cref{hier:comparison,hier:charges} for the first pass. For fixed
$M_0,C_0,n$, with $m,N$ subsequently large, this gives
\[
 \E[\one_{\mathcal P}K]
 \leq \frac{CD_*}{\log M_0}
       +o_{n\to\infty}(1)
       +o_{m,N\to\infty;n,M_0,C_0}(1).
\]
Choose $M_0$ large enough for the first term to meet the assigned
Markov bound at threshold $\kappa$, then take $n,m,N$ large in order.
Together with \cref{van:no-large-bound}, this proves
\cref{van:good-transmission} outside the assigned exception.

\emph{Internal jump energy.}
Inside $I$, the intermediate system has exactly the same data as the
full stage-$n$ system without artificial probes. Consider a deep square
of radius $R/4$ and its bounded-jump neighborhoods. A shield at fixed
relative radii inside $I$ has intact stage-$n$ tests. By
\cref{screen:projected,screen:paths}, comparison across that shield
changes the relevant scattering block by at most
$C_n e^{-c_nR^\alpha}$ in operator norm outside an event of vanishing
probability. The block ranks grow only polynomially in $R$.

Fix $n$ with $d_n\geq d/2$. By \cref{screen:limits}, choose a finite
jump cutoff capturing the spatial average on this deep square up to,
say, a quarter of $d_n$. For sufficiently large $m,N$, that truncated
average exceeds $d_n/2$ with the prescribed high probability. Every
such jump lies in $I$. Its squared chord length under $F$ is its
squared physical length divided by $R^2$, up to factors tending to
one as $b_N/R\to\infty$. The deep square has area comparable to
$R^2$. The preceding operator-norm comparison incurs a vanishing
error in the resulting weighted sum, because the rank is polynomial
and the shield error is exponential. This gives the lower bound in
\cref{van:good-jumps}, with $c_1>0$ independent of the later scales.

For the upper bound, extend $F|_I$ to data vanishing by radius $2R$,
with Lipschitz constant $C/R$. Apply the coarse profile bound in
\cref{screen:profiles} and then the first-pass comparison. The
expected intermediate energy is at most $CD_*+o(1)$. The sum of
internal jump terms is at most twice this energy. Markov's inequality
therefore gives a fixed upper cutoff $C_1$, chosen before $C_0$, for
the assigned exception budget.

\emph{Weighted transmission.}
For cell ports within a sufficiently large absolute multiple of
$L_*$ from $c$, chord bounds give $|F_\nu|\leq CM_0^2$. Their
contribution to $\norm{F_\nu T}_{\mathrm{HS}}$ is at most
$CM_0^2\sqrt K$. Beyond this range, retain only the good-edge costs
on the thick grid walls $O_D$. Those costs are still available in the
intermediate system, including their buffered tests. The wall estimate
in \cref{screen:walls}, followed by
\cref{screen:projected}, gives for a dyadic distance range of order
$D_1$ the expectation bound
\begin{equation}\label{van:far-weighted}
 C_{n,M_0}R^2(D_1/R)^2
 \left(e^{-c_nD_1^\alpha}+D_1^{-30}\right).
\end{equation}
To see the rank factor, there are $O_{M_0}(R^2)$ input probes; the
operator norm of transmission across a successful wall shield is
exponentially small, and it is always at most one by unitarity. The
chord weight is at most $CD_1/R$. At torus-size distances one uses a
separating band with radius a sufficiently small fixed multiple of
$D_1$, so the same estimates apply. The dyadic sum of
\cref{van:far-weighted} tends to zero as $m$ increases at fixed $n$,
uniformly for sufficiently large $N$: its polynomial part is a sum
of $D_1^{-28}$. Thus one may choose, for example,
$C_F=CM_0^2\sqrt\kappa+1$, with a sufficiently large absolute
constant and harmless extra slack for all four components. This
choice is independent of $n,C_0$; only the scale required to make the
far error small depends on them.

Finally include $\mathcal B$ using \cref{van:barrier}. The total
exception probability is less than $1/2$. This proves the event
assertion. Apply \cref{van:barrier-bound} with $D'=1$ and with
$D'=F_\nu|_J$. Since $S_{CJ}=(T^\top)_{CJ}$,
its unweighted Hilbert-Schmidt norm is at most $\sqrt K$, and its
weighted norm is at most $\norm{F_\nu T}_{\mathrm{HS}}$.
Substituting \cref{van:probe-scale} proves
\cref{van:external-small}.
\end{proof}

\subsection{An analytic trace with nonzero central value}

Close all inserted probes with the common phase $z$. For real $s$,
the corresponding physical shifts and the phase are
\begin{equation}\label{van:twist-parameter}
 v_x\longmapsto v_x+s\vartheta(v_x)\quad (x\in C),
 \qquad z=\frac{s+i\tau}{s-i\tau}.
\end{equation}
By \cref{mat:closure}, the scattering matrix on cell ports is
\begin{equation}\label{van:closed-family}
 S_z=A+\Delta(z),\qquad
 \Delta(z)=Tz(1-zB)^{-1}T^\top.
\end{equation}
Define $H(z)$ on the unit circle by replacing $S'$ with $S_z$ in
\cref{van:trace-data}. Its matrix expression contains the adjoint
$S_z^*$. To remove that obstruction to analyticity, replace the adjoint
by the adjugate and multiply the constant term by the determinant:
\begin{equation}\label{van:analytic-function}
 \mathfrak h(z)=\sum_\nu\tr\!\left[
 F_{I,\nu}\{F_\nu\det S_z-S_zF_\nu\operatorname{adj}S_z\}\right]
\end{equation}
This function is analytic in the disk. Indeed $B$ is a contraction, and the
scattering entries in \cref{van:closed-family} are rational analytic
functions there. Possible boundary poles are removable: away from
finitely many points their boundary values are unitary, which excludes
a genuine pole of a rational entry on the circle. Thus
$\mathfrak h$ extends continuously to the boundary. There,
\begin{equation}\label{van:boundary-modulus}
 \mathfrak h(z)=\det(S_z)H(z),\qquad
 |\mathfrak h(z)|=|H(z)|,
\end{equation}
because $\operatorname{adj}S_z=(\det S_z)S_z^*$.

\begin{lemma}\label{van:analytic-bounds}
Choose $\kappa$ sufficiently small depending on $d$, then the
constants in \cref{van:good}, and then $C_0$ sufficiently large.
There are constants $c_2>0$ and $B_p<\infty$, independent of the
later choices $n,m,N$, such that on $\mathcal G$
\begin{equation}\label{van:central-bound}
 |\mathfrak h(0)|\geq c_2,
\end{equation}
and, for every fixed $0<p<1/2$,
\begin{equation}\label{van:hardy-bound}
 \int_{\T}|H(z)|^p\,\dd\gamma_{\T}(z)\leq B_p.
\end{equation}
Here $\gamma_{\T}$ is normalized Haar measure. Neither constant
depends on the number of cell ports or probes.
\end{lemma}

\begin{proof}
Unitarity in \cref{van:block} gives $AA^*=1-TT^*$. If the
eigenvalues of $TT^*$ are $\delta_j\geq0$, then
$\sum_j\delta_j=K\leq\kappa$. Thus $A$ is invertible and
\begin{equation}\label{van:almost-unitary}
 |\det A|^2=\prod_j(1-\delta_j)\geq1-K,
 \qquad
 \norm{A^{-1}-A^*}_1
 =\norm{A^{-1}TT^*}_1
 \leq\frac{K}{\sqrt{1-K}}\leq2K
\end{equation}
for sufficiently small $\kappa$. At zero,
\begin{equation}\label{van:central-trace}
 \frac{\mathfrak h(0)}{\det A}
 =\sum_\nu\tr\!\left[F_{I,\nu}
                      (F_\nu-AF_\nu A^{-1})\right].
\end{equation}
On $I$ every $F_\nu$ is bounded by an absolute constant, because
the radius of $I$ is $R+O(b_m)$.

First retain only the $I,I$ indices in the product term of
\cref{van:central-trace}, along with the entire constant trace.
Replacing the internal block of $A^{-1}$ by that of $A^*$ costs at
most $CK$ in trace norm, by \cref{van:almost-unitary}. Symmetry of
$A$ makes the resulting internal expression exactly
\begin{equation}\label{van:internal-positive}
 \frac12\sum_{i,j\in I}|A_{ij}|^2|F(i)-F(j)|^2
 +\sum_{i\in I}|F(i)|^2
                  \left(1-\sum_{j\in I}|A_{ij}|^2\right).
\end{equation}
The second term is nonnegative, since $A$ is a contraction. The
first is at least $c_1d/2$ on $\mathcal G$.

For the external part, put
\[
 \delta_A=\norm{A_{IJ}}_{\mathrm{HS}}
       +\sum_\nu\norm{A_{IJ}F_\nu|_J}_{\mathrm{HS}}.
\]
Symmetry and \cref{van:almost-unitary} give
\[
 \norm{(A^{-1})_{JI}}_{\mathrm{HS}}
 \leq\norm{A_{IJ}}_{\mathrm{HS}}+2K.
\]
The external trace is consequently bounded in modulus by
$C\delta_A(\delta_A+K)$. Choose $\kappa$ so that the internal
replacement error is less than $c_1d/8$. After the choices in
\cref{van:good}, choose $C_0$ so large that
\cref{van:external-small} makes the external error less than
$c_1d/8$. The real part of \cref{van:central-trace} is then at
least $c_1d/4$. Combining with $|\det A|\geq\sqrt{1-\kappa}$
proves \cref{van:central-bound}.

For the boundary moment, expand the trace using $S_z=A+\Delta$.
The part $H_A$ with $A,A^*$ has a bounded modulus on $\mathcal G$.
Indeed its internal part is \cref{van:internal-positive}, whose
first term is bounded by $C_1/2$, and
\[
 \sum_{i\in I}\left(1-\sum_{j\in I}|A_{ij}|^2\right)
 \leq K+\norm{A_{IJ}}_{\mathrm{HS}}^2.
\]
The external part is bounded using
\cref{van:external-small}. All these bounds involve only the fixed
cutoffs in \cref{van:good}.

Consider a linear term, for one component,
\[
 \tr(F_I\Delta F A^*)
 =\tr\!\left[z(1-zB)^{-1}D_F\right],
 \qquad D_F=T^\top F A^*F_I T.
\]
Its coefficient satisfies
\[
 \norm{D_F}_1
 \leq\norm{T^\top F A^*}_{\mathrm{HS}}
       \norm{F_I T}_{\mathrm{HS}}
 \leq\norm{FT}_{\mathrm{HS}}
       \norm{F_I T}_{\mathrm{HS}}.
\]
The other linear term is its complex conjugate on the circle.
By \cref{mat:weak}, these terms have weak $L^1$ bounds depending
only on $C_F$ and $\kappa$, since
$\norm{F_I T}_{\mathrm{HS}}\leq C\sqrt\kappa$.

The quadratic term is bounded by
\[
 |\tr(F_I\Delta F\Delta^*)|
 \leq\norm{F_I\Delta}_{\mathrm{HS}}
       \norm{\Delta F}_{\mathrm{HS}}.
\]
Again \cref{mat:weak}, with the respective sandwiches
$(F_I T,T^\top)$ and $(T,T^\top F)$, gives weak $L^1$ bounds for
each factor. A nonnegative random variable with weak $L^1$ bound
$a$ has every moment $q<1$ bounded by $C_q a^q$, by integration
of its tail. Applying Cauchy--Schwarz with exponent $2p<1$ therefore
bounds the $p$-moment of this product. There are four components and
finitely many terms. Subadditivity at exponent $p<1$ proves
\cref{van:hardy-bound}. The use of Hilbert-Schmidt products and
trace norms makes the bounds independent of all matrix dimensions.
\end{proof}

\subsection{Changing the potential law under a bounded twist}

The artificial profile vanishes at the support endpoints and belongs
to $C^{1,1}([-h,h])$. For sufficiently small $|s|$, the map
\[
 T_s(v)=v+s\vartheta(v)
\]
is therefore a diffeomorphism of $[-h,h]$ onto itself, with derivative
at least $1/2$. If $r_s$ is the density of its pushforward of the
uniform law relative to that same law, then
\[
 r_s(y)=\frac{1}{1+s\vartheta'(T_s^{-1}(y))},
 \qquad |r_s-1|\leq C|s|,\qquad \int r_s\,\dd\mu=1.
\]
Consequently
\begin{equation}\label{van:single-density}
 \int r_s^2\,\dd\mu
 =1+\int(r_s-1)^2\,\dd\mu\leq1+Cs^2.
\end{equation}
Here $\mu$ is the uniform law on $[-h,h]$.

Fix $W<\infty$ and assume $|s|/\tau\leq W$. Conditional on the
offset and the test tile, change every true variable in $C$ by $T_s$
and leave all other true variables, selectors, and auxiliary uniforms
unchanged. There are at most $C(M_0R)^2$ changed sites. By independence
and \cref{van:single-density}, the resulting joint law has a density
$\mathscr D_s$ relative to the original joint law with
\begin{equation}\label{van:joint-density}
 \norm{\mathscr D_s}_{L^2}^2
 \leq \exp(CM_0^2R^2s^2)
 \leq \exp(CM_0^2C_0^2W^2)=D_W^2.
\end{equation}
This also holds after averaging over the offset and the test tile.
The bound is uniform in $m,N$ once they are sufficiently large for
the maps to be diffeomorphisms. Its constant may be large; only its
independence of the later scales matters.

We next check why this change of true potentials reproduces the
closed-probe system despite the potential-dependent hierarchy.
This check concerns recomputed statuses, not an assumption that
statuses are fixed by definition.

On $\mathcal P$, selected tiles of level at least $m$ do not meet
the test tile. Every tile of level below $m$ touching $C$ lies wholly
in the thinned gap, since $b_m\ll R$ and the annulus is separated
from both retained regions by a distance comparable to $R$. Thus a
selected tile capable of carrying an intermediate cell port is
disjoint from $C$. Its own true variables, own corner labels, local
likelihood, and reference potential are unchanged by the shift.
Eligibility depends only on selectors and is unchanged as well.

There is one remaining possibility: a changed coordinate may alter
the validity of an ancestor label used by a retained tile. Any such
ancestor contains both the retained root and a coordinate in $C$.
Their separation is at least a fixed multiple of $R$, so the ancestor
side is at least $cR$. The probability that any such large tile has
an invalid own label is bounded by
\begin{equation}\label{van:gate-exception}
 u_R=C_{M_0}\sum_{b_l\geq cR}(1+R/b_l)^2b_l^{-K_0}=o(1).
\end{equation}
This counts only tiles meeting $C$, which also covers every relevant
changed triple. The original bound is uniform in the offset and
torus size. After the shift the probability is at most
$D_W\sqrt{u_R}$ by \cref{van:joint-density}. Exclude both events.
All potentially changed ancestor validity tests are then valid before
and after the shift. Ancestor box identities themselves need not be
the same: the acceptance rules use those identities only through the
validity gates, whereas each retained tile's own box and likelihood
are unchanged.

On the remaining part of $\mathcal P$, all intermediate cell ports
and references are therefore unchanged when the hierarchy is
recomputed with the same selectors and auxiliary uniforms. Changes
of acceptance inside the gap have no effect, because those ports
have already been thinned, and $\mathcal P$ excludes the larger
ones. Before closure the real potentials on $C$ were the true ones;
the only added terms there were the artificial absorbers. Hence
\cref{mat:closure,van:twist-parameter} imply that $H(z)$ is exactly
the no-extra-probe intermediate trace $H^0$ evaluated at the shifted
data. Almost-everywhere inverse statements suffice throughout.

Let $a_{n,m,N}=\E[\one_{\mathcal P}|H^0|]$ in the original law.
Since $\mathcal P$ uses only unchanged selectors, Markov's inequality
and \cref{van:joint-density} give, for every fixed $\delta>0$,
\begin{equation}\label{van:twist-probability}
 \PP\{\mathcal G,\ |H(z)|>\delta\}
 \leq e_R(W)+D_W\left(\frac{a_{n,m,N}}{\delta}\right)^{1/2},
 \qquad |s|/\tau\leq W,
\end{equation}
where one may take $e_R(W)=u_R+D_W\sqrt{u_R}$. This tends to zero
as $m$ grows at fixed $W,M_0,C_0$. The estimate is uniform over
bounded twists; it does not assert simultaneous invariance for all
twists on a single realization. Only the probability bound for each
twist will be integrated.

\subsection{The contradiction}

\begin{proof}[Proof of \cref{van:zero}]
Retain the assumption $d>0$. Choose the constants in the following
order. First fix the barrier constants $d_0,\gamma$, then choose
$\kappa$ small enough for \cref{van:analytic-bounds}. Choose $M_0$
and the cutoffs $C_1,C_F$ using \cref{van:good}. Then choose $C_0$
sufficiently large for the same analytic lemma. All these choices
precede $n,m,N$. They provide a fixed central lower bound $c_2>0$
and a fixed boundary moment bound $B_p$ on an event $\mathcal G$
of probability at least $1/2$ once those scales are sufficiently large.

Fix $0<p'<p<1/2$. Under Haar measure, $s/\tau$ in
\cref{van:twist-parameter} has the standard Cauchy law. Denote its
tail probability beyond $W$ by $\beta_W$, so $\beta_W\to0$.
Apply H\"older's inequality on the product of the underlying
probability space and the circle. The pointwise good-data bound
\cref{van:hardy-bound} gives
\begin{align}
 &\E\left[\one_{\mathcal G}
        \int_\T |H(z)|^{p'}\,\dd\gamma_\T(z)\right]
 \label{van:small-boundary}\\
 &\quad\leq \delta^{p'}
   +B_p^{p'/p}\left[
       \beta_W^{1-p'/p}
       +\left\{e_R(W)+D_W
          \left(\frac{a_{n,m,N}}{\delta}\right)^{1/2}
        \right\}^{1-p'/p}\right].\notag
\end{align}
For completeness, the first term covers $|H|\leq\delta$. The
Haar tail costs the first H\"older term. On the bounded-twist part,
\cref{van:twist-probability} bounds the product measure of
$\{\mathcal G,|H|>\delta\}$, and the same $p$-moment bound gives
the second H\"older term. Tonelli's theorem applies to these
nonnegative quantities; no realization-wise uniformity in the twist
is used.

The right side can be made smaller than $c_2^{p'}/4$. To see the
order without interchanging limits, first choose $W$ so the tail term
is small, then choose $\delta$ so $\delta^{p'}$ is small. With these
fixed, $D_W$ is a fixed finite constant. Choose the tolerance in
\cref{van:adjustment} sufficiently small relative to $D_W$ and
$\delta$. Take $n$ large for that adjustment and the good-data
requirements; then take $m$ large for all fixed-$n$ spatial estimates,
the probe-strength constraint, and $e_R(W)$; finally take $N$ large
for the finite-torus approximations at this fixed pair. Increasing
$m,N$ as required preserves the preceding assertions. Thus
\cref{van:small-boundary} is indeed less than $c_2^{p'}/4$.

On the other hand, $|\mathfrak h|^{p'}$ is subharmonic in the disk.
For example, one may first use the nonnegative Laplacian of
$(|\mathfrak h|^2+\epsilon)^{p'/2}$ and then let
$\epsilon\downarrow0$. The continuous boundary extension and
\cref{van:boundary-modulus,van:central-bound} give on every good
realization
\[
 \int_\T|H(z)|^{p'}\,\dd\gamma_\T(z)
 =\int_\T|\mathfrak h(z)|^{p'}\,\dd\gamma_\T(z)
 \geq |\mathfrak h(0)|^{p'}\geq c_2^{p'}.
\]
Taking the expectation over $\mathcal G$ gives a lower bound
$c_2^{p'}/2$, contradicting the preceding upper bound. Therefore
$d=0$.
\end{proof}
